\documentclass[12pt,a4paper]{article}

\usepackage[utf8]{inputenc}
\usepackage[T1]{fontenc}
\usepackage[spanish]{babel}
\usepackage{geometry}
\usepackage{setspace}
\usepackage{float}
\usepackage{array}
\usepackage{booktabs}
\usepackage{graphicx}  % para \resizebox
\geometry{margin=2.5cm}
\setstretch{1.3}

\begin{document}

% ============================================================
% INTRODUCCION
% ============================================================

\section{Introducción}

El presente trabajo documenta la investigación comparativa y la
implementación práctica de herramientas de automatización de pruebas
y de análisis estático de código, aplicadas al proyecto integrador
del equipo: \textbf{Cal.com}, una plataforma open-source para la
programación de citas y gestión de disponibilidad, construida sobre
Next.js, TypeScript y Prisma, con PostgreSQL como motor de base de
datos. El repositorio oficial sobre el cual se trabaja se encuentra
disponible en \url{https://github.com/UP240080/EyMDSW-Calcom}.

A lo largo del cuatrimestre el equipo ha analizado tres
funcionalidades críticas de la plataforma --gestión de
disponibilidad, flujo de reserva de citas y panel administrativo--
bajo los atributos de calidad de usabilidad, compatibilidad y
accesibilidad. La presente tarea complementa ese análisis al
incorporar dos frentes adicionales: por un lado, una investigación
sobre el estado del arte de los frameworks de pruebas end-to-end
(E2E) y de las herramientas de análisis estático de código; por otro
lado, la aplicación práctica de Selenium con pytest y de SonarQube
directamente sobre el código del proyecto.

Se investigarán al menos tres frameworks de automatización E2E
(Selenium, Cypress y Playwright) y al menos tres herramientas de
análisis estático (SonarQube, ESLint y Semgrep), comparando sus
características técnicas para justificar su pertinencia frente al
stack tecnológico de Cal.com. Posteriormente, el equipo implementará
una suite de pruebas E2E con Page Object Model sobre un módulo real
del flujo de reserva, y ejecutará un análisis estático con
SonarQube para identificar issues de calidad, los cuales serán
priorizados en un plan de remediación de deuda técnica. Finalmente,
se documentará la integración de ambas prácticas en un pipeline de
integración continua.
\clearpage


% ============================================================
% DESARROLLO
% ============================================================

\section{Desarrollo}

Esta sección integra los seis componentes solicitados por la tarea:
investigación de frameworks E2E, investigación de herramientas de
análisis estático, implementación de pruebas con Selenium,
análisis estático con SonarQube, integración en CI y un plan de
remediación de deuda técnica. Las dos primeras subsecciones
(investigación) sientan las bases conceptuales que justifican las
decisiones tomadas en las subsecciones prácticas que les siguen.

% ============================================================
% 4.1
% ============================================================

\subsection{Investigación: frameworks de automatización E2E}

Las pruebas end-to-end (E2E) validan el comportamiento de una
aplicación desde la perspectiva del usuario final, simulando
interacciones reales sobre la interfaz en un navegador. Jorgensen
(2022) ubica este tipo de pruebas en el nivel más alto de la
pirámide de pruebas, ya que ejercitan el sistema integrado de punta
a punta, a diferencia de las pruebas unitarias que aíslan
componentes individuales. Para el caso de Cal.com, cuyo flujo de
reserva de citas involucra múltiples pasos secuenciales
(selección de tipo de evento, elección de fecha y hora, captura de
datos y confirmación), este tipo de prueba resulta indispensable
para validar la experiencia completa del usuario.

Se compararon tres de los frameworks más utilizados actualmente en
la industria: \textbf{Selenium}, \textbf{Cypress} y
\textbf{Playwright}.

Selenium es el framework más longevo de los tres y opera mediante el
protocolo WebDriver, lo que le permite comunicarse con prácticamente
cualquier navegador a través de su respectivo driver (García, 2022).
Su madurez se traduce en una comunidad extensa y en soporte para
múltiples lenguajes de programación, pero también en una curva de
aprendizaje más pronunciada, ya que el manejo de esperas y la
gestión de drivers recae casi por completo en quien escribe la
prueba (García, 2022).

Cypress, en cambio, se ejecuta dentro del mismo bucle de eventos del
navegador, lo que le da acceso directo al DOM y a las peticiones de
red de la aplicación bajo prueba. Esta arquitectura simplifica
considerablemente el manejo de esperas, ya que Cypress reintenta
automáticamente las aserciones hasta que se cumplen o se agota un
tiempo límite (Aniche, 2022). Su principal limitación es que está
acotado al ecosistema JavaScript/TypeScript y, históricamente, a
navegadores basados en Chromium, Firefox y Edge.

Playwright, desarrollado por Microsoft, combina varias de las
ventajas de los dos anteriores: ofrece soporte multinavegador real
(Chromium, Firefox y WebKit) desde un único motor de automatización,
esperas automáticas integradas (auto-waiting) y soporte para
múltiples lenguajes, incluyendo TypeScript, Python, Java y C#. Su
arquitectura basada en contextos de navegador aislados le permite
ejecutar pruebas en paralelo con mayor velocidad que las
arquitecturas tradicionales basadas en WebDriver (Khorikov, 2020,
respecto a la importancia de la velocidad de ejecución para
mantener suites de prueba sostenibles).

\begin{table}[h]
\centering
\resizebox{\textwidth}{!}{%
\begin{tabular}{lccccccc}
\toprule
\textbf{Criterio} & \textbf{Selenium} & \textbf{Cypress} & \textbf{Playwright} \\
\midrule
Lenguajes soportados & Java, Python, C\#, JS/TS, Ruby, Kotlin & JavaScript / TypeScript & JS/TS, Python, Java, C\# \\
Curva de aprendizaje & Alta (gestión manual de drivers y esperas) & Media (API declarativa, propia de JS) & Media-baja (auto-wait, API moderna) \\
Velocidad de ejecución & Media (overhead del protocolo WebDriver) & Alta (ejecución dentro del navegador) & Alta (contextos aislados, ejecución paralela nativa) \\
Manejo de esperas & Manual (\texttt{WebDriverWait}) & Automático (retry-ability integrado) & Automático (auto-waiting integrado) \\
Soporte de navegadores & Chrome, Firefox, Edge, Safari (vía drivers) & Chrome, Edge, Firefox (Safari experimental) & Chromium, Firefox y WebKit nativos \\
Comunidad & Muy amplia, desde 2004 & Amplia, en crecimiento desde 2017 & En crecimiento acelerado, respaldo de Microsoft \\
Licencia & Apache 2.0 (open source) & MIT (núcleo open source; Cypress Cloud es de pago) & Apache 2.0 (open source) \\
\bottomrule
\end{tabular}%
}
\caption{Comparativa de frameworks de automatización E2E. Fuente: elaboración propia con base en García (2022), Aniche (2022) y documentación oficial (selenium.dev, docs.cypress.io, playwright.dev).}
\label{tab:e2e}
\end{table}

\textbf{Recomendación.} Dado que Cal.com es una aplicación construida
íntegramente en TypeScript sobre Next.js, Playwright sería la
opción más adecuada para un entorno de producción: ofrece soporte
nativo del lenguaje del proyecto, cobertura multinavegador real y
una velocidad de ejecución superior, factores relevantes para una
aplicación con flujos críticos como la reserva de citas. Sin
embargo, y tal como lo permite el lineamiento de la tarea, la
implementación práctica de este trabajo (sección 4.3) se realiza con
\textbf{Selenium + pytest}, por ser la herramienta trabajada en las
sesiones de clase y por contar con un ecosistema robusto de
documentación y patrones (como Page Object Model) ampliamente
documentados por García (2022). Esta decisión no contradice la
recomendación anterior, sino que distingue entre la herramienta
técnicamente más idónea para el stack del proyecto y la herramienta
pedagógicamente adecuada para los objetivos de esta unidad.


% ============================================================
% 4.2
% ============================================================

\subsection{Investigación: análisis estático de código}

El análisis estático de código examina el código fuente sin
ejecutarlo, identificando patrones que indican posibles defectos,
vulnerabilidades de seguridad, código duplicado o violaciones a
buenas prácticas de diseño (Khorikov, 2020). Esto lo distingue del
análisis dinámico --como las pruebas E2E descritas en la sección
anterior-- que observa el comportamiento real del sistema durante su
ejecución y, por lo tanto, solo puede detectar defectos en las rutas
de código que efectivamente se ejercitan durante la prueba. El
análisis estático, en cambio, recorre el cien por ciento del código
fuente, lo que le permite detectar problemas de forma temprana, antes
de que el sistema esté siquiera en ejecución, aunque con el costo de
generar en ocasiones falsos positivos.

Se compararon tres herramientas de análisis estático relevantes para
el stack de Cal.com: \textbf{SonarQube}, \textbf{ESLint} y
\textbf{Semgrep}.

SonarQube es una plataforma integral que evalúa cinco dimensiones de
calidad (confiabilidad, seguridad, mantenibilidad, cobertura y
duplicación) y soporta de forma nativa más de 30 lenguajes,
incluyendo TypeScript y JavaScript. Su fortaleza principal es ofrecer
una visión consolidada del estado de calidad de todo un repositorio,
junto con un mecanismo de \emph{Quality Gate} que puede bloquear
integraciones que no cumplan los umbrales definidos.

ESLint, por su parte, es un linter especializado únicamente en
JavaScript y TypeScript, enfocado en detectar errores de sintaxis,
problemas de estilo y antipatrones específicos del lenguaje en
tiempo de escritura, generalmente integrado directamente en el editor
o como parte de los hooks de pre-commit. A diferencia de SonarQube,
no ofrece una visión histórica ni un tablero centralizado de
calidad, pero su retroalimentación es prácticamente inmediata para
quien programa.

Semgrep se distingue por su enfoque basado en reglas semánticas
ligeras, definidas en patrones legibles similares al propio código,
lo que facilita escribir reglas personalizadas para detectar
vulnerabilidades de seguridad específicas (por ejemplo, uso indebido
de consultas SQL sin parametrizar) de forma más ágil que con motores
de análisis más pesados (Aniche, 2022, respecto a la importancia de
detectar errores comunes de forma temprana y automatizada).

\begin{table}[h]
\centering
\resizebox{\textwidth}{!}{%
\begin{tabular}{lccc}
\toprule
\textbf{Criterio} & \textbf{SonarQube} & \textbf{ESLint} & \textbf{Semgrep} \\
\midrule
Tipo de defectos detectados & Bugs, vulnerabilidades, code smells, duplicación, hotspots de seguridad & Errores de sintaxis, antipatrones de JS/TS, estilo de código & Vulnerabilidades de seguridad y patrones personalizados \\
Lenguajes soportados & +30 lenguajes (incluye TS/JS) & JavaScript / TypeScript (vía plugins) & +20 lenguajes, reglas personalizables \\
Integración en flujo de trabajo & Servidor centralizado / CI, Quality Gate & Editor / pre-commit / CI & CLI / CI, reglas como código \\
Visión histórica del proyecto & Sí (tablero con tendencias) & No & Limitada \\
Licencia & Community Edition gratuita; ediciones avanzadas de pago & MIT (open source) & Open source (núcleo); Semgrep Cloud de pago \\
\bottomrule
\end{tabular}%
}
\caption{Comparativa de herramientas de análisis estático de código. Fuente: elaboración propia con base en Khorikov (2020), Aniche (2022) y documentación oficial (docs.sonarsource.com, eslint.org, semgrep.dev).}
\label{tab:static}
\end{table}

\textbf{Justificación para el stack de Cal.com.} SonarQube resulta
adecuado para este proyecto porque analiza de forma nativa
TypeScript y JavaScript --los lenguajes en los que está escrito
prácticamente todo el código de Cal.com--, y porque su mecanismo de
\emph{Quality Gate} permite establecer un criterio objetivo y
automatizable sobre el cual el equipo puede decidir si un cambio es
apto para integrarse al proyecto, algo especialmente valioso dado el
volumen y la complejidad del repositorio descritos como riesgo en la
elección del caso de estudio. Además, SonarQube puede consumir
directamente el archivo \texttt{coverage.xml} generado por
\texttt{pytest-cov}, lo que permite vincular en un mismo tablero la
cobertura producida por la suite de pruebas E2E (sección 4.3) con el
resto de las métricas estáticas del código (sección 4.4). Por estas
razones, se selecciona SonarQube como herramienta principal de
análisis estático para la práctica, complementando su uso, de ser
necesario, con las reglas de ESLint ya presentes en la configuración
original del proyecto Cal.com.

% ============================================================
% Referencias usadas en esta parte (para fusionar en la
% sección 7 de Referencias del documento final):
%
% Aniche, M. (2022). Effective Software Testing: A Developer's
%   Guide. Manning.
% García, B. (2022). Hands-On Selenium WebDriver with Java.
%   O'Reilly.
% Jorgensen, P. C. (2022). Software Testing: A Craftsman's
%   Approach (5a ed.). CRC Press.
% Khorikov, V. (2020). Unit Testing Principles, Practices, and
%   Patterns. Manning.
% Selenium. (s.f.). Selenium Documentation.
%   https://www.selenium.dev/documentation
% Cypress. (s.f.). Cypress Documentation.
%   https://docs.cypress.io
% Playwright. (s.f.). Playwright Documentation.
%   https://playwright.dev
% SonarSource. (s.f.). SonarQube Documentation.
%   https://docs.sonarsource.com
% ESLint. (s.f.). ESLint Documentation. https://eslint.org
% ============================================================


\end{document}